%--------------------------------------------------------------------------------------
% ÁLTALÁNOS ALAPSABLON A (RÉSZ)BEADANDÓKHOZ
%--------------------------------------------------------------------------------------
\documentclass[12pt, a4paper]{book}
\usepackage[margin=1in]{geometry}
\usepackage[utf8]{inputenc}
\usepackage[magyar]{babel}
\usepackage[hidelinks]{hyperref} 
\usepackage{graphicx}
\usepackage{comment}
\usepackage{amsthm}
\usepackage{MnSymbol}
\usepackage{wasysym}
\usepackage{fancyhdr}
\usepackage{longtable}
\usepackage{changepage}
\usepackage{xcolor}
\usepackage{algorithmicx}
\usepackage{listings}

\pagestyle{fancy}
\fancyhf{}
\fancyhead[LE,RO]{\thepage}
\fancyhead[RE,LO]{\leftmark}
%kitöltendő
\fancyfoot[CE,CO]{Csapatsorszám - Csapatnév}

\title{Szoftver projekt labor beadandó} 
%kitöltendő
\author{Karsa Zoltán, Killin Attila, Jakab Dániel, Kacsó Péter, Imets Tamás} 

\setcounter{secnumdepth}{4}

\begin{document}
\begin{titlepage} 
	\begin{center}
		\includegraphics[width=5cm]{doc/img/BMElogo}\\ 
		\vspace{1cm}
		\LARGE{\bfseries{Szoftver projekt laboratórium}}\\
		\vspace{0.5cm}
		%kitöltendő
		\Large\textsc{X. Összesített dokumentum}\\
		\vspace{1cm}
		\small{Csapat}\\
		%kitöltendő
		\Large{\bfseries{Csapatsorszám - Csapatnév}}\\
		\vspace{1cm}
		\small{Konzulens}\\
		%kitöltendő
		\Large{Dr. Goldschmidt Balázs}
	\end{center}
	\begin{flushleft}
		\vspace*{6cm}
		Csapattagok\\
		\vspace{-0.3cm}
		\rule{14cm}{0.5pt}\\
		\vspace{0.2cm}
		\begin{tabular}{l l l} %kitöltendő
			Teszt Tamás & ?????? & teszttamas@iit.bme.hu \\
			Java Dániel & ?????? & javadaniel@iit.bme.hu \\
			Projlab Péter & ?????? & projlabpeter@iit.bme.hu \\
			LaTex Zoltán  & ?????? & latexzoltan@iit.bme.hu \\
			Git Attila & ?????? & gitattila@iit.bme.hu \\
		\end{tabular}
	\end{flushleft}
	\begin{flushright}
		\vspace*{2cm}
		\today   %vagy saját dátum megadása, manuálisan
	\end{flushright}
\end{titlepage}
%---------------------------------------------------------------
%\\\\\\\\\\\\\\\\\\\\\\\\NE MÓDOSÍTSD///////////////////////////

\newenvironment{naplo}{\begin{longtable}{|l|l|p{2.5cm}|p{7.5cm}|}
        \hline
        \hline \multicolumn{1}{|c|}{\textbf{Kezdet}} & \multicolumn{1}{c|}{\textbf{Időtartam}} & \multicolumn{1}{c|}{\textbf{Résztvevők}} & \multicolumn{1}{c|}{\textbf{Leírás}}  \\ \hline 
        \endfirsthead
        
        \multicolumn{3}{c}%
        {{\bfseries}} \\
        \hline \multicolumn{1}{|c|}{\textbf{Kezdet}} & \multicolumn{1}{c|}{\textbf{Időtartam}} & \multicolumn{1}{c|}{\textbf{Résztvevők}} & \multicolumn{1}{c|}{\textbf{Leírás}}  \\ \hline 
        \endhead
        
}{\end{longtable}}

\newcommand{\naplotag}[4]{#1 & #2 & #3 & #4 \\
    \hline}

\newtheorem{defi}{Definíció}
\newtheorem{rov}{Rövidítés}
\newcommand{\urlref}[2]{
	\parbox{\linewidth}{\href{#1}{\RIGHTarrow \quad #2 - \texttt{#1}}}}

\newcommand{\diagram}[3]{
	\begin{figure}[h]
		\centering
		\includegraphics[width=#3]{#1}
		\caption{#2}
	\end{figure}}

\renewcommand{\comment}[1]{\begin{flushright}
		\color{red} 
		#1
\end{flushright} }

\newenvironment{class-template}[5]
{ 
    \hspace*{0.4em} $\blacksquare $ \large Felelősség \\  \normalsize #1 \\[12px]
    \hspace*{0.4em} $\blacksquare $ \large Ősosztályok \\  \normalsize #2 \\[12px]
    \hspace*{0.4em} $\blacksquare $ \large Interfészek \\  \normalsize #3 \\[12px]
    \hspace*{0.4em} $\blacksquare $ \large Attribútumok \\  \normalsize #4 \\[12px]
    \hspace*{0.4em} $\blacksquare $ \large Asszociációk \\  \normalsize #4 \\[12px]
    \hspace*{0.4em} $\blacksquare $ \large Metódusok \\  \normalsize #5 \\[12px]
}{}

\newenvironment{class-template-responsibility}{\hspace*{0.4em} $\blacksquare $ \large Felelősség \\[5px]  \normalsize}{\\[12px]}
\newenvironment{class-template-baseclass}{\hspace*{0.4em} $\blacksquare $ \large Ősosztály \\[5px]  \normalsize}{\\[12px]}
\newenvironment{class-template-interface}{\hspace*{0.4em} $\blacksquare $ \large Interfészek \\[5px]  \normalsize}{\\[12px]}
\newenvironment{class-template-method}{\hspace*{0.4em} $\blacksquare $ \large Metódusok \\[5px]  \normalsize}{\\[12px]}
\newenvironment{class-template-attribute}{\hspace*{0.4em} $\blacksquare $ \large Attribútumok \\[5px]  \normalsize}{\\[12px]}
\newenvironment{class-template-association}{\hspace*{0.4em} $\blacksquare $ \large Asszociációk \\[5px]  \normalsize}{\\[12px]}
\newenvironment{class-template-statechart}{\hspace*{0.4em} $\blacksquare$ \large Állapotgép \\[5px] \normalsize}{\\[12px]}

\newenvironment{enum-template}[2]
{ 
    \hspace*{0.4em} $\blacksquare $ \large Felelősség \\[5px]  \normalsize #1 \\[12px]
    \hspace*{0.4em} $\blacksquare $ \large Literálok \\[5px]  \normalsize #2 \\[12px]
}{}	

{\vspace*{-1cm}}
\newcommand{\classitem}[2]{
    \hspace*{25px}\parbox{0.9\linewidth}{$\Diamond $ \textit{#1} - #2}
}

\newcommand{\baseclass}{ $\rightarrow$ }

\newenvironment{funkovetelmeny}[6]
{
    \noindent\begin{tabular}{| p{2cm} | p{2cm} | p{2.2cm} | p{5.3cm} | p{2.5cm} |}
        \hline
        \hline
        \textbf{Azonosító} & \textbf{Prioritás} & \textbf{Forrás} & \textbf{Use-case} & \textbf{Ellenőrzés} \\
        \hline
        #1 & #2 & #3 & #4 & #5 \\
        \hline
    \end{tabular}
    \\[1ex]
    \textbf{Leírás:} #6 \\[0.2ex]
    \itshape}
{\\[5px]}

\newenvironment{kovetelmeny}[5]
{ 
    \noindent\begin{tabular}{| p{2cm} | p{3.5cm} | p{3.5cm} | p{3.5cm} |}
        \hline
        \hline
        \textbf{Azonosító} & \textbf{Prioritás} & \textbf{Forrás} & \textbf{Ellenőrzés} \\
        \hline
        #1 & #2 & #3 & #4 \\
        \hline
    \end{tabular}
    \\[1ex]
    \textbf{Leírás:} #5 \\[0.2ex]
    \itshape}
{\\[5px]}

\newenvironment{use-case}[3]
{ 
    
    \vspace{10px}
    \noindent\begin{tabular}{| p{3cm} | p{12cm} |}
        \hline
        \hline
        \vspace{0.2ex} \textsc{\textbf{Use-case neve:}} & \vspace{0.2ex} \textsc{#1} \\[0.2ex]
        \hline
        \textbf{Rövid leírás:} & #2 \\
        \hline
        \textbf{Aktorok:} & #3 \\
        \hline
        \textbf{Forgatókönyv:} &  
    }
    {	\\
        \hline
\end{tabular} }

\newenvironment{fajllista}{
    \begin{longtable}{|p{4.5cm}|l|l| p{4.5cm} |}
        \hline
        \hline \multicolumn{1}{|c|}{\textbf{Fájl neve}} & \multicolumn{1}{c|}{\textbf{Méret}} & \multicolumn{1}{c|}{\textbf{Keletkezés ideje}} & \multicolumn{1}{c|}{\textbf{Tartalom}}  \\ \hline 
        \endfirsthead
        
        \multicolumn{3}{c}%
        {{\bfseries}} \\
        \hline \multicolumn{1}{|c|}{\textbf{Fájl neve}} & \multicolumn{1}{c|}{\textbf{Méret}} & \multicolumn{1}{c|}{\textbf{Keletkezés ideje}} & \multicolumn{1}{c|}{\textbf{Tartalom}}  \\ \hline 
        \endhead
        
    }{\end{longtable}}

\newcommand{\fajl}[4] {
    #1 & #2 & #3 & #4 \\ \hline
}

\newenvironment{ertekeles}{
    \begin{figure}[h]
        \centering
        \noindent\begin{tabular}{| p{5cm} | p{4cm} | p{4cm} |}
            \hline
            \hline
            \textbf{Tag neve} & \textbf{Tag neptun} & \textbf{Munka százalékban} \\	\hline
}{\end{tabular}\end{figure}}

\newcommand{\ertekelestag}[3]{ 
    #1 & #2 & #3 \\ \hline
}

\newcommand{\ertekelestagk}[2]{ 
    #1 & #2 \\ \hline
}

\newenvironment{ertekelesOra}{
    \begin{figure}[h]
        \centering
        \noindent\begin{tabular}{| p{5cm} | p{4cm} | p{4cm} |}
            \hline
            \hline
            \textbf{Tag neve} & \textbf{Tag neptun} & \textbf{Munka órában} \\	\hline
            
        }{\end{tabular}	\end{figure}}

\newenvironment{ertekelesKod}{
    \begin{figure}[h]
        \centering
        \noindent\begin{tabular}{| p{5cm} | p{4cm} |}
            \hline
            \hline
            \textbf{Fázis} & \textbf{Kódsorok száma} \\	\hline
            
        }{\end{tabular}	\end{figure}}

\newcommand{\test}[3]{
    \vspace{10px}
    \noindent\begin{tabular}{| p{3.5cm} | p{12cm} |}
        \hline
        \hline
        \vspace{0.2ex} \textsc{\textbf{Teszt-eset neve}} & \vspace{0.2ex} \textsc{#1} \\[0.2ex]
        \hline
        \textbf{Rövid leírás} & #2 \\
        \hline
        \textbf{Teszt célja} & #3 \\
        \hline
    \end{tabular}
}

\newenvironment{test-case}[4]
{ 
    \hspace*{0.4em} $\blacksquare $ \large Leírás   \\  \normalsize #1   \\[12px]
    \hspace*{0.4em} $\blacksquare $ \large Ellenőrzött funkcionalitás, várható hibahelyek  \\  \normalsize #2  \\[12px]
    \hspace*{0.4em} $\blacksquare $ \large Bemenet  \\  \normalsize #3  \\[12px]
    \hspace*{0.4em} $\blacksquare $ \large Elvárt kimenet \\  \normalsize #4 \\[12px]
}{}

\newenvironment{test-case-description}{\hspace*{0.4em} $\blacksquare $ \large Leírás \\[5px]  \normalsize}{\\[12px]}
\newenvironment{test-case-function}{\hspace*{0.4em} $\blacksquare $ \large Ellenőrzött funkcionalitás, várható hibahelyek \\[5px]  \normalsize}{\\[12px]}
\newenvironment{test-case-input}{\hspace*{0.4em}\parbox{\linewidth}{$\blacksquare $ \large Bemenet\vspace*{5px}}\normalsize}{\leavevmode\newline}
\newenvironment{test-case-output}{\hspace*{0.4em}\parbox{\linewidth}{$\blacksquare $ \large Elvárt kimenet\vspace{5px}}\normalsize}{\leavevmode\newline}

\newcommand{\testOK}[2]{
    \vspace{10px}
    \noindent\begin{tabular}{| p{4.2cm} | p{12cm} |}
        \hline
        \hline
        \textbf{Tesztelő neve} & \textsc{#1} \\[0.2ex]
        \hline
        \textbf{Teszt időpontja} & #2 \\
        \hline
    \end{tabular}
}

\newcommand{\testFAIL}[5]{
    \vspace{10px}
    \noindent\begin{tabular}{| p{4.2cm} | p{12cm} |}
        \hline
        \hline
        \textbf{Tesztelő neve} & \textsc{#1} \\[0.2ex]
        \hline
        \textbf{Teszt időpontja} & #2 \\
        \hline
        \textbf{Teszt eredménye} & #3 \\
        \hline
        \textbf{Lehetséges hibaokok} & #4 \\
        \hline
        \textbf{Változtatások} & #5 \\
        \hline
    \end{tabular}
}

\newenvironment{terv}{	
    \begin{longtable}{| l | p{8cm} | l |}
        \hline
        \hline
        \multicolumn{1}{|c|}{\textbf{Határidő}} & \multicolumn{1}{c|}{\textbf{Feladat}} & \multicolumn{1}{c|}{\textbf{Felelős}}  \\ \hline 
        \endfirsthead
        
        \multicolumn{3}{c}%
        {{\bfseries}} \\
        \hline \multicolumn{1}{|c|}{\textbf{Határidő}} & \multicolumn{1}{c|}{\textbf{Feladat}} & \multicolumn{1}{c|}{\textbf{Felelős}}  \\ \hline 
        \endhead
        
}{\end{longtable}}

\newcommand{\tervitem}[3]{ #1 & #2 & #3 \\ \hline}

\newenvironment{szotar}{
    \begin{longtable}{ p{4cm}  p{11cm} }
        
}{\end{longtable}}

\newcommand{\szotaritem}[2]{ \textbf{#1} & #2 \\[10px]}

\sloppy

%-----------------------------------------------------------------------------------------

%Végső beadáshoz tartalomjegyzék:
%\tableofcontents

%A fejezetek számozásának átállítására (mindig a megadott számtól folytatódik a számozás):
%\setcounter{chapter}{-1}


\setcounter{chapter}{1}
\chapter{Követelmény, projekt, funkcionalitás}

\section{Bevezetés}

\subsection{Cél}
\comment{A dokumentum célja.}

\subsection{Szakterület}
\comment{A kialakítandó szoftver milyen területen használható, milyen célra.}

\subsection{Definíciók, rövidítések}
\comment{A dokumentumban használt definíciók, rövidítések magyarázata.}
\begin{defi} 
	Példa defi: egy demó a defi használatára
\end{defi}
\begin{rov}
    PR: példa rövidítés
\end{rov}

\subsection{Hivatkozások}
\urlref{http://iit.bme.hu/}{BME IIT - Programozás alapjai 3. segédanyagok, Szoftvertechnikák segédanyagok, Szoftver projekt laboratórium feladatok}

\subsection{Összefoglalás}
\comment{A dokumentum további részeinek rövid ismertetése}

\section{Áttekintés}
\subsection{Általános áttekintés}
\comment{A kialakítandó szoftver legmagasabb szintű architekturális képe. A fontosabb alrendszerek felsorolása, a közöttük kialakítandó interfészek lényege, a felhasználói kapcsolatok alapja. Esetleges hálózati és adattárolási elvárások.}

\subsection{Funkciók}
\comment{A feladat kb. 4000 karakteres (kb 1,5 oldal) részletezettségű magyar nyelvű leírása. Nem szerepelhetnek informatikai kifejezések.}

\subsection{Felhasználók}
\comment{A felhasználók jellemzői, tulajdonságai}

\subsection{Korlátozások}
\comment{Az elkészítendő szoftverre vonatkozó – általában nem funkcionális - előírások, korlátozások.}

\subsection{Feltételezések, kapcsolatok}
\comment{A dokumentumban használt anyagok, web-oldalak felsorolása}


\section{Követelmények}
\subsection{Funkcionális követelmények}
\comment{Az alábbi táblázat kitöltésével készítendő. Dolgozzon ki követelmény azonosító rendszert! Az ellenőrzés módja szokásosan bemutatás és/vagy kiértékelés. Prioritás lehet alapvető, fontos, opcionális. Az alapvető követelmények nem teljesítése végzetes. Forrás alatt a követelményt előíró anyagot, szervezetet kell érteni. Esetünkben forrás lehet maga a csapat is, mikor ő talál ki követelményt. Use-case-ek alatt az adott követelményt megvalósító használati esete(ke)t kell megadni.}
\begin{funkovetelmeny}
	%Azonosító
	{Azonosító}
	%Prioritás
	{Pioritás}
	%Forrás
	{Forrás}
	%Use case
	{Használati eset}
	%Ellenőrzés
	{llenőrzés}
	%Leírás
	{Leírás}
    Megjegyzés (opcionális)
\end{funkovetelmeny}


\subsection{Erőforrásokkal kapcsolatos követelmények}
\comment{A szoftver fejlesztésével és használatával kapcsolatos számítógépes, hardveres, alapszoftveres és egyéb architekturális és logisztikai követelmények}

\begin{kovetelmeny}
	{Azonosító} %Azonosító
	{Pioritás} %Prioritás
	{Forrás} %Forrás
	{Ellenőrzés} %Ellenőrzés
	{Leírás} %Leírás
    Megjegyzés (opcionális)
\end{kovetelmeny}

\subsection{Átadással kapcsolatos követelmények}
\comment{A szoftver átadásával, telepítésével, üzembe helyezésével kapcsolatos követelmények}
\begin{kovetelmeny}
    {Azonosító} %Azonosító
    {Pioritás} %Prioritás
    {Forrás} %Forrás
    {Ellenőrzés} %Ellenőrzés
    {Leírás} %Leírás
    Megjegyzés (opcionális)
\end{kovetelmeny}

\subsection{Egyéb nem funkcionális követelmények}
\comment{A biztonsággal, hordozhatósággal, megbízhatósággal, tesztelhetőséggel, a felhasználóval kapcsolatos követelmények}
\begin{kovetelmeny}
    {Azonosító} %Azonosító
    {Pioritás} %Prioritás
    {Forrás} %Forrás
    {Ellenőrzés} %Ellenőrzés
    {Leírás} %Leírás
    Megjegyzés (opcionális)
\end{kovetelmeny}

\section{Lényeges use-case-ek}
\subsection{Use-case leírások}

\begin{use-case}
	%név
	{Use-case Neve}
	%rövid leírás
	{Az eset rövid leírása}
	%aktorok
	{Aktorok}
	%forgatókönyv
	Forgatókönyv \newline 
        \textbf{A.1} Alternatíva
\end{use-case}

\subsection{Use-case diagram}
\diagram{doc/img/BMElogo}{Demó}{4cm}

\section{Szótár}
\comment{A szótár a követelmények alapján készítendő fejezet. Egy szótári bejegyzés definiálásához csak más szótári bejegyzések és köznapi – a feladattól független – fogalmak használhatók fel. A szótár mérete kb. 1-2 oldal legyen.}

%szotar.sh script, lsd. a dokumentációt!!!
\begin{szotar}
    \szotaritem{Kulcs 1 }{Érték 1}
    \szotaritem{Hóásás }{Egy réteg hó eltakarítása, azaz törlése az adott jégtábláról. Egy egység munkába kerül.}
\end{szotar}

\section{Projekt terv}
\comment{Tartalmaznia kell a projekt végrehajtásának lépéseit, a lépések, eredmények határidejét, az egyes feladatok elvégzéséért felelős személyek nevét és beosztását, a szükséges erőforrásokat, stb. Meg kell adni a csoportmunkát támogató eszközöket, a választott technikákat! Definiálni kell, hogy hogyan történik a dokumentumok és a forráskód megosztása!}


\subsection{Projektütemterv}
\begin{terv}
		\tervitem{febr. 4.}  {Követelmény, projekt, funkcionalitás}{ ?????? }
        \tervitem{márc. 2. } { Analízis modell kidolgozása 1. - beadás }{ ?????? } 
        \tervitem{márc. 9. } { Analízis modell kidolgozása 2. - beadás }{ ?????? }
        \tervitem{márc. 16. }{ Skeleton tervezése - beadás }{ ?????? }
        \tervitem{márc. 23. }{ Skeleton - beadás és a forráskód herculesre való feltöltése }{ ?????? }
        \tervitem{márc. 30. }{ Prototípus koncepciója - beadás }{ ?????? }
        \tervitem{ápr. 6. }{ Részletes tervek - beadás }{ ?????? }
        \tervitem{ápr. 27. }{ Prototípus - beadás és a forráskód, a tesztbemenetek és az elvárt kimenetek herculesre való feltöltése }{ ?????? }
        \tervitem{ máj. 4. }{ Grafikus felület specifikációja - beadás }{ ?????? }
        \tervitem{ máj. 18. }{ Grafikus változat és Összefoglalás - beadás és a forráskód herculesre való feltöltése }{ ?????? } 
\end{terv}

\subsection{Erőforrások, eszközök}
A fejlesztés során felhasznált segédeszközök:
\begin{itemize}
	\item Dokumentáció: dokumentáló eszközök felsorlása
	\item Kommunikáció: kommunikációs platformok...
	\item Modellező eszköz: modellezési eszközök
	\item Fejlesztő környezetek: ...
	\item Forráskód megosztás, verziókezelés: ...
	\item Egyéb ....
\end{itemize}


\clearpage
\section{Napló}

\comment{ 
    A napló tartalmazza az előző beadás óta eltelt időszak történéseit időrendben. A naplóból egyértelműen ki kell derülnie, hogy az egyes anyagrészeket ki és mennyi idő alatt készítette.
    
    A napló bejegyzésekből áll. Minden bejegyzésnek tartalmaznia kell: 
    \begin{itemize}
        \item a történés kezdetének időpontját, nap-óra pontossággal
        \item történés időtartamát, óra felbontással
        \item a szereplő(k) nevét (Kérjük a szereplők VEZETÉKNEVÉT használni)
        \item a tevékenység leírását.
    \end{itemize}

    Amennyiben a tevékenységben több szereplő vesz részt, akkor aza tevékenység csak értekezlet lehet, amelynek az eredményei DÖNTÉSEK. A döntéseket precízen meg kell szövegezni (Pl.: Az X objektum Y és Z metódusainak kódját W készíti el Q határidőre).Ha a bejegyzés egyetlen személyhez kötődik, akkor meg kell adni, hogy a tevékenység milyen dologra irányul. A dolog a feladat kapcsán elkészítendő termék, amelynek a (esetleg korábban) beadott anyagban megtalálhatónak kell lenni.A naplóbejegyzés felbontásának egysége szöveges, rajzos anyag esetében az ábra, diagram, vagy kb. fél-egy oldalnyi szöveg. Kódban az egység a metódus. (Pl.: A 3. ábrán látható szekvencia-diagram kidolgozása, vagy az X objektum Y és Z metódusainak kódolása és belövése.
}

\begin{naplo}
	\naplotag{feb. 18. 16h }{ 1 óra }{ Csapat }{ Értekezlet  \newline Döntés: Segédeszközök kiválasztása (git, trello, drive)}
\end{naplo}

\chapter{Analízis modell kidolgozása 1.}

\section{Objektum katalógus}

\comment{Minden, a feladatban szereplő objektum rövid, egy-két bekezdés hosszú ismertetése. Meg kell jelenjen minden objektumhoz, hogy mi a felelőssége. Informális leírás, ezért nem kell foglalkozni az örökléssel, az interfészekkel, az absztrakt osztályokkal, a segédosztályokkal.}

\subsection{Objektum1}
\comment{Felelősség informális leírása}

\subsection{Objektum2}
\comment{Felelősség informális leírása}

\comment{Az objektumkatalógus alapján kiindulva kell megalkotni az objektumorientált analízis modellt. A 3.2, 3.3, 3.4, 3.5 alfejezetek ugyanannak a modellnek a különböző nézetei, ezért egyidőben, egymással összefüggésben készülnek.Megtörténik az objektumkatalógusbantárgyalt objektumok felelősségének formalizálásaosztályokká,attribútumokká, metódusokká. Csak publikus metódusok szerepelhetnek. Megjelennek az interfészek, az öröklés, az absztrakt osztályok. Segédosztályokra még nincs szükség.}

\section{Statikus struktúra diagramok}
\comment{Az előző alfejezet osztályainak kapcsolatait és publikus metódusait bemutató osztálydiagram(ok). Tipikus hibalehetőségek: csillag-topológia, szigetek.}

\diagram{doc/img/BMElogo}{Demó}{3cm}

\section{Osztályok leírása}
\comment{Az előző alfejezetben tárgyalt objektumok felelősségének formalizálása attribútumokká, metódusokká. Csak publikus metódusok szerepelhetnek. Ebben az alfejezetben megjelennek az interfészek, az öröklés, az absztrakt osztályok. Segédosztályokra még mindig nincs szükség. Az osztályok ABC sorrendben kövessék egymást. Interfészek esetén az Interfészek, Attribútumok pontok kimaradnak.}

\subsection{Osztály1}
\begin{class-template-responsibility}
    Felelősség leírása
\end{class-template-responsibility}
\begin{class-template-interface}
    Megvalósított interfészek felsorolása
\end{class-template-interface}
\begin{class-template-baseclass}
    Ős-Ősosztály \baseclass Ősosztály... 
\end{class-template-baseclass}
\begin{class-template-attribute}
    \classitem{+A [0..*]}{adattag A}
\end{class-template-attribute}
\comment{Milyen publikus, protected és privát  metódusokkal rendelkezik. Metódusonként egy-három mondat arról, hogy mit csinál.}
\begin{class-template-method}
    \classitem{+B(A a) : void}{metódus B}
\end{class-template-method}

\subsection{Osztály2}
\begin{class-template-responsibility}
    Felelősség leírása
\end{class-template-responsibility}
\begin{class-template-statechart}
    A belső működéshez tartozó állapotgép
    \diagram{doc/img/BMElogo}{Az osztály állapotváltozásai}{3cm}
\end{class-template-statechart}

\section{Statikus struktúra diagramok}
\comment{Az előző alfejezet osztályainak kapcsolatait és publikus metódusait bemutató osztálydiagram(ok). Tipikus hibalehetőségek: csillag-topológia, szigetek.}
\diagram{doc/img/BMElogo}{osztálydiagram}{3cm}


\section{Szekvencia diagramok}
\comment{Inicializálásra, use-case-ekre, belső működésre. Konzisztens kell legyen az előző alfejezettel. Minden metódus, ami ott szerepel, fel kell tűnjön valamelyik szekvenciában. Minden metódusnak, ami szekvenciában szerepel, szereplnie kell a valamelyik osztálydiagramon.}
\diagram{doc/img/BMElogo}{Szekvencia1}{3cm}

\section{State-chartok}
\comment{Csak azokhoz az osztályokhoz, ahol van értelme. Egyetlen állapotból álló state-chartok ne szerepeljenek. A játék működését bemutató state-chart-ot készíteni tilos.}

\clearpage
\section{Napló}

\comment{ 
    A napló tartalmazza az előző beadás óta eltelt időszak történéseit időrendben. A naplóból egyértelműen ki kell derülnie, hogy az egyes anyagrészeket ki és mennyi idő alatt készítette.
    
    A napló bejegyzésekből áll. Minden bejegyzésnek tartalmaznia kell: 
    \begin{itemize}
        \item a történés kezdetének időpontját, nap-óra pontossággal
        \item történés időtartamát, óra felbontással
        \item a szereplő(k) nevét (Kérjük a szereplők VEZETÉKNEVÉT használni)
        \item a tevékenység leírását.
    \end{itemize}
    
    Amennyiben a tevékenységben több szereplő vesz részt, akkor aza tevékenység csak értekezlet lehet, amelynek az eredményei DÖNTÉSEK. A döntéseket precízen meg kell szövegezni (Pl.: Az X objektum Y és Z metódusainak kódját W készíti el Q határidőre).Ha a bejegyzés egyetlen személyhez kötődik, akkor meg kell adni, hogy a tevékenység milyen dologra irányul. A dolog a feladat kapcsán elkészítendő termék, amelynek a (esetleg korábban) beadott anyagban megtalálhatónak kell lenni.A naplóbejegyzés felbontásának egysége szöveges, rajzos anyag esetében az ábra, diagram, vagy kb. fél-egy oldalnyi szöveg. Kódban az egység a metódus. (Pl.: A 3. ábrán látható szekvencia-diagram kidolgozása, vagy az X objektum Y és Z metódusainak kódolása és belövése.
}

\begin{naplo}
    \naplotag{feb. 18. 16h }{ 1 óra }{ Csapat }{ Értekezlet  \newline Döntés: Segédeszközök kiválasztása (git, trello, drive)}
\end{naplo}

\chapter{Analízis modell kidolgozása 2.}

\section{Objektum katalógus}

\comment{Minden, a feladatban szereplő objektum rövid, egy-két bekezdés hosszú ismertetése. Meg kell jelenjen minden objektumhoz, hogy mi a felelőssége. Informális leírás, ezért nem kell foglalkozni az örökléssel, az interfészekkel, az absztrakt osztályokkal, a segédosztályokkal.}

\subsection{Objektum1}
\comment{Felelősség informális leírása}

\subsection{Objektum2}
\comment{Felelősség informális leírása}

\comment{Az objektumkatalógus alapján kiindulva kell megalkotni az objektumorientált analízis modellt. A 3.2, 3.3, 3.4, 3.5 alfejezetek ugyanannak a modellnek a különböző nézetei, ezért egyidőben, egymással összefüggésben készülnek.Megtörténik az objektumkatalógusbantárgyalt objektumok felelősségének formalizálásaosztályokká,attribútumokká, metódusokká. Csak publikus metódusok szerepelhetnek. Megjelennek az interfészek, az öröklés, az absztrakt osztályok. Segédosztályokra még nincs szükség.}

\section{Statikus struktúra diagramok}
\comment{Az előző alfejezet osztályainak kapcsolatait és publikus metódusait bemutató osztálydiagram(ok). Tipikus hibalehetőségek: csillag-topológia, szigetek.}

\diagram{doc/img/BMElogo}{Demó}{3cm}

\section{Osztályok leírása}
\comment{Az előző alfejezetben tárgyalt objektumok felelősségének formalizálása attribútumokká, metódusokká. Csak publikus metódusok szerepelhetnek. Ebben az alfejezetben megjelennek az interfészek, az öröklés, az absztrakt osztályok. Segédosztályokra még mindig nincs szükség. Az osztályok ABC sorrendben kövessék egymást. Interfészek esetén az Interfészek, Attribútumok pontok kimaradnak.}

\subsection{Osztály1}
\begin{class-template-responsibility}
    Felelősség leírása
\end{class-template-responsibility}
\begin{class-template-interface}
    Megvalósított interfészek felsorolása
\end{class-template-interface}
\begin{class-template-baseclass}
    Ős-Ősosztály \baseclass Ősosztály... 
\end{class-template-baseclass}
\begin{class-template-attribute}
    \classitem{+A [0..*]}{adattag A}
\end{class-template-attribute}
\comment{Milyen publikus, protected és privát  metódusokkal rendelkezik. Metódusonként egy-három mondat arról, hogy mit csinál.}
\begin{class-template-method}
    \classitem{+B(A a) : void}{metódus B}
\end{class-template-method}

\subsection{Osztály2}
\begin{class-template-responsibility}
    Felelősség leírása
\end{class-template-responsibility}
\begin{class-template-statechart}
    A belső működéshez tartozó állapotgép
    \diagram{doc/img/BMElogo}{Az osztály állapotváltozásai}{3cm}
\end{class-template-statechart}

\section{Statikus struktúra diagramok}
\comment{Az előző alfejezet osztályainak kapcsolatait és publikus metódusait bemutató osztálydiagram(ok). Tipikus hibalehetőségek: csillag-topológia, szigetek.}
\diagram{doc/img/BMElogo}{osztálydiagram}{3cm}


\section{Szekvencia diagramok}
\comment{Inicializálásra, use-case-ekre, belső működésre. Konzisztens kell legyen az előző alfejezettel. Minden metódus, ami ott szerepel, fel kell tűnjön valamelyik szekvenciában. Minden metódusnak, ami szekvenciában szerepel, szereplnie kell a valamelyik osztálydiagramon.}
\diagram{doc/img/BMElogo}{Szekvencia1}{3cm}

\section{State-chartok}
\comment{Csak azokhoz az osztályokhoz, ahol van értelme. Egyetlen állapotból álló state-chartok ne szerepeljenek. A játék működését bemutató state-chart-ot készíteni tilos.}

\clearpage
\section{Napló}

\comment{ 
    A napló tartalmazza az előző beadás óta eltelt időszak történéseit időrendben. A naplóból egyértelműen ki kell derülnie, hogy az egyes anyagrészeket ki és mennyi idő alatt készítette.
    
    A napló bejegyzésekből áll. Minden bejegyzésnek tartalmaznia kell: 
    \begin{itemize}
        \item a történés kezdetének időpontját, nap-óra pontossággal
        \item történés időtartamát, óra felbontással
        \item a szereplő(k) nevét (Kérjük a szereplők VEZETÉKNEVÉT használni)
        \item a tevékenység leírását.
    \end{itemize}
    
    Amennyiben a tevékenységben több szereplő vesz részt, akkor aza tevékenység csak értekezlet lehet, amelynek az eredményei DÖNTÉSEK. A döntéseket precízen meg kell szövegezni (Pl.: Az X objektum Y és Z metódusainak kódját W készíti el Q határidőre).Ha a bejegyzés egyetlen személyhez kötődik, akkor meg kell adni, hogy a tevékenység milyen dologra irányul. A dolog a feladat kapcsán elkészítendő termék, amelynek a (esetleg korábban) beadott anyagban megtalálhatónak kell lenni.A naplóbejegyzés felbontásának egysége szöveges, rajzos anyag esetében az ábra, diagram, vagy kb. fél-egy oldalnyi szöveg. Kódban az egység a metódus. (Pl.: A 3. ábrán látható szekvencia-diagram kidolgozása, vagy az X objektum Y és Z metódusainak kódolása és belövése.
}

\begin{naplo}
    \naplotag{feb. 18. 16h }{ 1 óra }{ Csapat }{ Értekezlet  \newline Döntés: Segédeszközök kiválasztása (git, trello, drive)}
\end{naplo}

\chapter{Szkeleton tervezése}


\section{A szkeleton modell valóságos use-case-ei}
\comment{A szkeletonnak, mint önálló programnak a működésével kapcsolatos use-case-ek.}

\subsection{Use-case diagram}
\diagram{doc/img/BMElogo}{Demó}{4cm}

\subsection{Use-case leírások}
\comment{Minden use-case-hez külön}
\begin{use-case}
    %név
    {Use-case Neve}
    %rövid leírás
    {Az eset rövid leírása}
    %aktorok
    {Aktorok}
    %forgatókönyv
    Forgatókönyv \newline 
    \textbf{A.1} Alternatíva
\end{use-case}

\section{A szkeleton kezelői felületének terve, dialógusok}
\comment{A szkeleton által elfogadott bemenetek , valamint a szöveges konzolon megjelenő kimenetek. A kiemenet formátuma olyan kell legyen, ami alapján a működés összevethető a korábbi szekvencia-diagramokkal.}

\comment{Konzol input/output szemléltetésre:}
\begin{verbatim}
    move 1 to 3
    get item
\end{verbatim}

\section{Szekvencia diagramok a belső működésre}
\comment{A szkeletonban implementált szekvenciadiagramok. Tipikusan egy use-case egy diagram. Ezek megegyezhetnek a korábban specifikált diagramokkal, de az egyes életvonalakat (lifeline) egyértelműen a szkeletonban példányosított objektumokhoz kell tudni kötni. Azt kell megjeleníteni, hogy a szkeletonban létrehozott objektumok egymással hogyan fognak kommunikálni.}

\section{Kommunikációs diagramok}
\comment{A szkeletonban, az egyes szkeleton-use-case-ek futása során létrehozott objektumok és kapcsolataik bemutatására szolgáló diagramok. Ezek alapján valósítják meg a szkeleton fejlesztői az inicializáló kódrészleteket.}

\clearpage
\section{Napló}

\comment{ 
    A napló tartalmazza az előző beadás óta eltelt időszak történéseit időrendben. A naplóból egyértelműen ki kell derülnie, hogy az egyes anyagrészeket ki és mennyi idő alatt készítette.
    
    A napló bejegyzésekből áll. Minden bejegyzésnek tartalmaznia kell: 
    \begin{itemize}
        \item a történés kezdetének időpontját, nap-óra pontossággal
        \item történés időtartamát, óra felbontással
        \item a szereplő(k) nevét (Kérjük a szereplők VEZETÉKNEVÉT használni)
        \item a tevékenység leírását.
    \end{itemize}
    
    Amennyiben a tevékenységben több szereplő vesz részt, akkor aza tevékenység csak értekezlet lehet, amelynek az eredményei DÖNTÉSEK. A döntéseket precízen meg kell szövegezni (Pl.: Az X objektum Y és Z metódusainak kódját W készíti el Q határidőre).Ha a bejegyzés egyetlen személyhez kötődik, akkor meg kell adni, hogy a tevékenység milyen dologra irányul. A dolog a feladat kapcsán elkészítendő termék, amelynek a (esetleg korábban) beadott anyagban megtalálhatónak kell lenni.A naplóbejegyzés felbontásának egysége szöveges, rajzos anyag esetében az ábra, diagram, vagy kb. fél-egy oldalnyi szöveg. Kódban az egység a metódus. (Pl.: A 3. ábrán látható szekvencia-diagram kidolgozása, vagy az X objektum Y és Z metódusainak kódolása és belövése.
}

\begin{naplo}
    \naplotag{feb. 18. 16h }{ 1 óra }{ Csapat }{ Értekezlet  \newline Döntés: Segédeszközök kiválasztása (git, trello, drive)}
\end{naplo}

\chapter{Szkeleton beadás}

\section{Fordítási és futtatási útmutató}
\comment{A feltöltött program fordításával és futtatásával kapcsolatos útmutatás. Ennek tartalmaznia kell leltárszerűen az egyes fájlok pontos nevét, méretét byte-ban, keletkezési idejét, valamint azt, hogy a fájlban mi került megvalósításra.}

\subsection{Fájllista}

\begin{fajllista}
	\fajl
	{Demo.java} % Kezdet
	{353 byte} % Méret
	{2020.03.26~21:05~} % Résztvevők
	{Demo programosztály} % Leírás
\end{fajllista}

\subsection{Fordítás}
\comment{A fenti listában szereplő forrásfájlokból milyen műveletekkel lehet a bináris, futtatható kódot előállítani. Az előállításhoz csak a 2. Követelmények c. dokumentumban leírt környezetet szabad előírni.}

\begin{verbatim}
    javac -d bin *.java
\end{verbatim}

\subsection{Futtatás}
\comment{A futtatható kód elindításával kapcsolatos teendők leírása. Az indításhoz csak a 2. Követelmények c. dokumentumban leírt környezetet szabad előírni.}

\begin{verbatim}
    cd bin
    java Main.java
\end{verbatim}

\clearpage
\section{Értékelés}
\comment{A projekt kezdete óta az értékelésig eltelt időben tagokra bontva, százalékban.}
\begin{ertekeles}
	  \ertekelestag{Teszt János}{ ?????? }{100\%}
\end{ertekeles}
\clearpage
\section{Napló}

\comment{ 
    A napló tartalmazza az előző beadás óta eltelt időszak történéseit időrendben. A naplóból egyértelműen ki kell derülnie, hogy az egyes anyagrészeket ki és mennyi idő alatt készítette.
    
    A napló bejegyzésekből áll. Minden bejegyzésnek tartalmaznia kell: 
    \begin{itemize}
        \item a történés kezdetének időpontját, nap-óra pontossággal
        \item történés időtartamát, óra felbontással
        \item a szereplő(k) nevét (Kérjük a szereplők VEZETÉKNEVÉT használni)
        \item a tevékenység leírását.
    \end{itemize}
    
    Amennyiben a tevékenységben több szereplő vesz részt, akkor aza tevékenység csak értekezlet lehet, amelynek az eredményei DÖNTÉSEK. A döntéseket precízen meg kell szövegezni (Pl.: Az X objektum Y és Z metódusainak kódját W készíti el Q határidőre).Ha a bejegyzés egyetlen személyhez kötődik, akkor meg kell adni, hogy a tevékenység milyen dologra irányul. A dolog a feladat kapcsán elkészítendő termék, amelynek a (esetleg korábban) beadott anyagban megtalálhatónak kell lenni.A naplóbejegyzés felbontásának egysége szöveges, rajzos anyag esetében az ábra, diagram, vagy kb. fél-egy oldalnyi szöveg. Kódban az egység a metódus. (Pl.: A 3. ábrán látható szekvencia-diagram kidolgozása, vagy az X objektum Y és Z metódusainak kódolása és belövése.
}

\begin{naplo}
    \naplotag{feb. 18. 16h }{ 1 óra }{ Csapat }{ Értekezlet  \newline Döntés: Segédeszközök kiválasztása (git, trello, drive)}
\end{naplo}

\chapter{Prototípus koncepciója}
\comment{\hrule}
\comment{Amennyiben változott a modell:}
\setcounter{section}{-1}
\section{Változás hatása a modellre}
\subsection{Módosult osztálydiagram}
\comment{Az analízis modell osztálydiagramja a változások figyelembevételével.}
\subsection{Új vagy megváltozott metódusok}
\comment{Az analízis modell osztályleírásaiból azon metódusok újbóli felsorolása leírással együtt, amelyek a változtatás miatt módosultak vagy újonnan be lettek vezetve}
\textbf{A bevezetett új osztályok}
\subsubsection{Osztály1}
\begin{class-template-responsibility}
    Felelősség leírása
\end{class-template-responsibility}
\begin{class-template-baseclass}
    Ős-Ősosztály \baseclass Ősosztály... 
\end{class-template-baseclass}
\begin{class-template-method}
    \classitem{+B(A a) : void}{metódus B}
\end{class-template-method}
\textbf{A korábbi osztályokban történt módosítások}
\subsubsection{Osztály2}
\begin{class-template-responsibility}
    Felelősség leírása
\end{class-template-responsibility}
\begin{class-template-baseclass}
    Ős-Ősosztály \baseclass Ősosztály... 
\end{class-template-baseclass}
\begin{class-template-method}
    \classitem{+B(A a) : void}{metódus B}
\end{class-template-method}

\subsection{Szekvencia-diagramok}
\comment{Az analízis modell szekvenciadiagramjaiból a változás által érintett, előírt, módosítottdiagramok}
\diagram{doc/img/BMElogo}{Megváltozott szekvenciák}{3cm}
\comment{\hrule}

\section{Prototípus interface-definíciója}
\comment{Definiálni kell a teszteket leíró nyelvet. Külön figyelmet kell fordítani arra, hogy ha a rendszer véletlen elemeket is tartalmaz, akkor a véletlenszerűség ki-bekapcsolható legyen, és a program determinisztikusan is tesztelhető legyen.}

\subsection{Az interfész általános leírása}
\comment{A protó (karakteres) input és output felületeit úgy kell kialakítani, hogy az input fájlból is vehető legyen illetőleg az output fájlba menthető legyen, vagyis kommunikációra csak a szabványos be- és kimenet használható.}

\subsection{Bemeneti nyelv}
\comment{Definiálni kell a teszteket leíró nyelvet. Külön figyelmet kell fordítani arra, hogy ha a rendszer véletlen elemeket is tartalmaz, akkor a véletlenszerűség ki-bekapcsolható legyen, és a program determinisztikusan is futtatható legyen. A szálkezelést is tesztelhető, irányítható módon kell megoldani. A programot egy adott konfigurációból is el kell tudni indítani, vagyis kell olyan parancs, amivel konkrét előre megadott állapotból indul a rendszer (pl. load).}

\begin{itemize}
    \item Parancs1
    \begin{itemize}
        \item Leírás:
        \item Opciók:
    \end{itemize}
    \item Parancs2
    \begin{itemize}
        \item Leírás:
        \item Opciók:
    \end{itemize}
\end{itemize}
Példák a bemeneti nyelvre:
\begin{verbatim}
    parancs1 arg1 arg2 arg3
    parancs2 arg1 arg2
\end{verbatim}

\comment{Ha szükséges, meg kell adni a konfigurációs (pl. pályaképet megadó) fájlok nyelvtanát is.}

\subsection{Kimeneti nyelv}
\comment{Egyértelműen definiálni kell, hogy az egyes bemeneti parancsok végrehajtása után előálló állapot milyen formában jelenik meg a szabványos kimeneten. A program képes legyen olyan kimenetet előállítani, amellyel az objektumok állapota ellenőrizhető (pl. save). Ebben az alfejezetben is precízen definiálni kell, hogy a kimenet nyelve milyen elemekből és milyen szintakszissal áll elő.}
\begin{verbatim}
    kimeneti szintaktika
\end{verbatim}

\section{Összes részletes use-case}
\begin{use-case}
    %név
    {Use-case Neve}
    %rövid leírás
    {Az eset rövid leírása}
    %aktorok
    {Aktorok}
    %forgatókönyv
    Forgatókönyv \newline 
    \textbf{A.1} Alternatíva
\end{use-case}

\section{Tesztelési terv}
\comment{A tesztelési tervben definiálni kell, hogy a be- és kimeneti fájlok egybevetésével miként végezhető el a program tesztelése. Meg kell adni teszt forgatókönyveket. Az egyes teszteket elég informálisan, szabad szövegként leírni. Teszt-esetenként egy-öt mondatban. Minden teszthez meg kell adni, hogy mi a célja, a proto mely funkcionalitását, osztályait stb. teszteli. Az alábbi táblázat minden teszt-esethez külön-külön elkészítendő.}

\test{Teszteset neve}{Leírás}{A végrehajtott teszt célja}

\section{Tesztelést támogató segéd- és fordítóprogramok specifikálása}
\comment{Specifikálni kell a tesztelést támogató segédprogramokat.}
\clearpage
\section{Napló}

\comment{ 
    A napló tartalmazza az előző beadás óta eltelt időszak történéseit időrendben. A naplóból egyértelműen ki kell derülnie, hogy az egyes anyagrészeket ki és mennyi idő alatt készítette.
    
    A napló bejegyzésekből áll. Minden bejegyzésnek tartalmaznia kell: 
    \begin{itemize}
        \item a történés kezdetének időpontját, nap-óra pontossággal
        \item történés időtartamát, óra felbontással
        \item a szereplő(k) nevét (Kérjük a szereplők VEZETÉKNEVÉT használni)
        \item a tevékenység leírását.
    \end{itemize}
    
    Amennyiben a tevékenységben több szereplő vesz részt, akkor aza tevékenység csak értekezlet lehet, amelynek az eredményei DÖNTÉSEK. A döntéseket precízen meg kell szövegezni (Pl.: Az X objektum Y és Z metódusainak kódját W készíti el Q határidőre).Ha a bejegyzés egyetlen személyhez kötődik, akkor meg kell adni, hogy a tevékenység milyen dologra irányul. A dolog a feladat kapcsán elkészítendő termék, amelynek a (esetleg korábban) beadott anyagban megtalálhatónak kell lenni.A naplóbejegyzés felbontásának egysége szöveges, rajzos anyag esetében az ábra, diagram, vagy kb. fél-egy oldalnyi szöveg. Kódban az egység a metódus. (Pl.: A 3. ábrán látható szekvencia-diagram kidolgozása, vagy az X objektum Y és Z metódusainak kódolása és belövése.
}

\begin{naplo}
    \naplotag{feb. 18. 16h }{ 1 óra }{ Csapat }{ Értekezlet  \newline Döntés: Segédeszközök kiválasztása (git, trello, drive)}
\end{naplo}

\chapter{Részletes tervek}
\comment{A dokumentum célja, hogy pontosan specifikálja az implementálandó osztályokat, beleértve a privát attribútumokat és metódusokat, ezek definícióját is.A dokumentum második fele részletesen be kell mutassa a korábban definiált be-és kimeneti nyelv szintakszisát felhasználva, hogy mely tesztekkel lesz a prototípus ellenőrizve.}
\section{Osztályok és metódusok tervei}
\subsection{Osztály1}
\begin{class-template-responsibility}
    Felelősség leírása
\end{class-template-responsibility}
\begin{class-template-interface}
    Megvalósított interfészek felsorolása
\end{class-template-interface}
\begin{class-template-baseclass}
    Ős-Ősosztály \baseclass Ősosztály... 
\end{class-template-baseclass}
\begin{class-template-attribute}
    \classitem{+A [0..*]}{adattag A}
\end{class-template-attribute}
\comment{Milyen publikus, protected és privát  metódusokkal rendelkezik. Metódusonként precíz leírás, ha szükséges, activity diagram is a metódusban megvalósítandó algoritmusról. Minden olyan metódusnak szerepelnie kell, amelyiket az osztály megvalósít vagy felüldefiniál.}
\begin{class-template-method}
    \classitem{+B(A a) : void}{metódus B}
\end{class-template-method}

\subsection{Osztály2}
\begin{class-template-responsibility}
    Felelősség leírása
\end{class-template-responsibility}
\begin{class-template-statechart}
    A belső működéshez tartozó állapotgép
    \diagram{doc/img/BMElogo}{Az osztály állapotváltozásai}{3cm}
\end{class-template-statechart}

\section{A tesztek részletes tervei, leírásuk a teszt nyelvén}
\comment{A tesztek részletes tervei alatt meg kell adni azokat a bemeneti adatsorozatokat, amelyekkel a program működése ellenőrizhető. Minden bemenő adatsorozathoz definiálni kell, hogy az adatsorozat végrehajtásától a program mely részeinek, funkcióinak ellenőrzését várjuk és konkrétan milyen eredményekre számítunk, ezek az eredmények hogyan vethetők össze a bemenetekkel.A tesztek leírásakor az előző dokumentumban (proto koncepciója) megadott szintakszist kell használni.}

\subsection{Teszteset1}
\begin{test-case-description}
    Szöveges leírás, kb 1-5 mondat
\end{test-case-description}
\begin{test-case-function}
    ...
\end{test-case-function}
\begin{test-case-input}
    A proto nyelvén megadva, pl.:
    \begin{verbatim}
    init world
    0: S 3
    ...
    print
    \end{verbatim}
\end{test-case-input}
\begin{test-case-output}
    A poro kimeneti nyelvén megadva, pl.:
    \begin{verbatim}
    worlddata
    ...
    creatures 0
    \end{verbatim}
\end{test-case-output}

\section{A tesztelést támogató programok tervei}
\comment{A tesztadatok előállítására, a tesztek eredményeinek kiértékelésére szolgáló segédprogramok részletes terveit kell elkészíteni.}
\clearpage
\section{Napló}

\comment{ 
    A napló tartalmazza az előző beadás óta eltelt időszak történéseit időrendben. A naplóból egyértelműen ki kell derülnie, hogy az egyes anyagrészeket ki és mennyi idő alatt készítette.
    
    A napló bejegyzésekből áll. Minden bejegyzésnek tartalmaznia kell: 
    \begin{itemize}
        \item a történés kezdetének időpontját, nap-óra pontossággal
        \item történés időtartamát, óra felbontással
        \item a szereplő(k) nevét (Kérjük a szereplők VEZETÉKNEVÉT használni)
        \item a tevékenység leírását.
    \end{itemize}
    
    Amennyiben a tevékenységben több szereplő vesz részt, akkor aza tevékenység csak értekezlet lehet, amelynek az eredményei DÖNTÉSEK. A döntéseket precízen meg kell szövegezni (Pl.: Az X objektum Y és Z metódusainak kódját W készíti el Q határidőre).Ha a bejegyzés egyetlen személyhez kötődik, akkor meg kell adni, hogy a tevékenység milyen dologra irányul. A dolog a feladat kapcsán elkészítendő termék, amelynek a (esetleg korábban) beadott anyagban megtalálhatónak kell lenni.A naplóbejegyzés felbontásának egysége szöveges, rajzos anyag esetében az ábra, diagram, vagy kb. fél-egy oldalnyi szöveg. Kódban az egység a metódus. (Pl.: A 3. ábrán látható szekvencia-diagram kidolgozása, vagy az X objektum Y és Z metódusainak kódolása és belövése.
}

\begin{naplo}
    \naplotag{feb. 18. 16h }{ 1 óra }{ Csapat }{ Értekezlet  \newline Döntés: Segédeszközök kiválasztása (git, trello, drive)}
\end{naplo}

\setcounter{chapter}{9}
\chapter{Prototípus beadása}

\section{Fordítási és futtatási útmutató}
\comment{A feltöltött program fordításával és futtatásával kapcsolatos útmutatás. Ennek tartalmaznia kell leltárszerűen az egyes fájlok pontos nevét, méretét byte-ban, keletkezési idejét, valamint azt, hogy a fájlban mi került megvalósításra.}

\subsection{Fájllista}

\begin{fajllista}
    \fajl
    {Demo.java} % Kezdet
    {353 byte} % Méret
    {2020.03.26~21:05~} % Időpont
    {Demo programosztály} % Leírás
\end{fajllista}

\subsection{Fordítás}
\comment{A fenti listában szereplő forrásfájlokból milyen műveletekkel lehet a bináris, futtatható kódot előállítani. Az előállításhoz csak a 2. Követelmények c. dokumentumban leírt környezetet szabad előírni.}

\begin{verbatim}
javac -d bin *.java
\end{verbatim}

\subsection{Futtatás}
\comment{A futtatható kód elindításával kapcsolatos teendők leírása. Az indításhoz csak a 2. Követelmények c. dokumentumban leírt környezetet szabad előírni.}

\begin{verbatim}
cd bin
java Main.java
\end{verbatim}

\section{Tesztek jegyzőkönyvei}

\subsection{Teszteset1}
\comment{Az alábbi táblázatot a megismételt (hibás) tesztek esetén kell kitölteni minden ismétléshez egyszer. Ha szükséges, akkor a valós kimenet is mellékelhető mint a teszt eredménye}
\testFAIL{????}{Ápr. 25. 13:50}{ A teszt eredménye}{A teszt hibás futásának okai}{Változtatások}
\comment{Az alábbi táblázatot az utolsó, sikeres tesztfuttatáshoz kell kitölteni}
\testOK{????}{Ápr. 25. 14:00}

\section{Értékelés}
\comment{A projekt kezdete óta az értékelésig eltelt időben tagokra bontva, százalékban.}
\begin{ertekeles}
    \ertekelestag{Teszt János}{ ?????? }{100\%}
\end{ertekeles}

\clearpage
\section{Napló}

\comment{ 
    A napló tartalmazza az előző beadás óta eltelt időszak történéseit időrendben. A naplóból egyértelműen ki kell derülnie, hogy az egyes anyagrészeket ki és mennyi idő alatt készítette.
    
    A napló bejegyzésekből áll. Minden bejegyzésnek tartalmaznia kell: 
    \begin{itemize}
        \item a történés kezdetének időpontját, nap-óra pontossággal
        \item történés időtartamát, óra felbontással
        \item a szereplő(k) nevét (Kérjük a szereplők VEZETÉKNEVÉT használni)
        \item a tevékenység leírását.
    \end{itemize}
    
    Amennyiben a tevékenységben több szereplő vesz részt, akkor aza tevékenység csak értekezlet lehet, amelynek az eredményei DÖNTÉSEK. A döntéseket precízen meg kell szövegezni (Pl.: Az X objektum Y és Z metódusainak kódját W készíti el Q határidőre).Ha a bejegyzés egyetlen személyhez kötődik, akkor meg kell adni, hogy a tevékenység milyen dologra irányul. A dolog a feladat kapcsán elkészítendő termék, amelynek a (esetleg korábban) beadott anyagban megtalálhatónak kell lenni.A naplóbejegyzés felbontásának egysége szöveges, rajzos anyag esetében az ábra, diagram, vagy kb. fél-egy oldalnyi szöveg. Kódban az egység a metódus. (Pl.: A 3. ábrán látható szekvencia-diagram kidolgozása, vagy az X objektum Y és Z metódusainak kódolása és belövése.
}

\begin{naplo}
    \naplotag{feb. 18. 16h }{ 1 óra }{ Csapat }{ Értekezlet  \newline Döntés: Segédeszközök kiválasztása (git, trello, drive)}
\end{naplo}

\chapter{Grafikus felület specifikációja}

\section{A grafikus interfész}
\comment{A menürendszer, a kezelői felület grafikus képe. A grafikus felület megjelenését, a használt ikonokat, stb screenshot-szerű képekkel kell bemutatni. Az építészetben ez a homlokzati terv.}
\diagram{doc/img/BMElogo}{Példa kép}{3cm}

\section{A grafikus rendszer architektúrája}
\comment{A felület működésének elve, a grafikus rendszer architektúrája (struktúra diagramok). A struktúra diagramokon a prototípus azon és csak azon osztályainak is szerepelnie kell, amelyekhez a grafikus felületet létrehozó osztályok kapcsolódnak.}

\subsection{A felület működési elve}
\comment{Le kell írni, hogy a grafikai megjelenésért felelős osztályok, objektumok hogyan kapcsolódnak a meglevő rendszerhez, a megjelenítés során mi volt az alapelv. Törekedni kell az MVC megvalósításra. Alapelvek lehetnek: \textbf{push} alapú: a modell értesíti a felületet, hogy változott; \textbf{pull} alapú: a felület kérdezi le a modellt, hogy változott-e; \textbf{kevert}: a kettő kombinációja.}

\subsection{A felület osztály-struktúrája}
\comment{Osztálydiagram. Minden új osztály, és azon régiek, akik az újakhoz közvetlenül kapcsolódnak.}
\diagram{doc/img/BMElogo}{Példa kép}{3cm}

\section{A grafikus objektumok felsorolása}
\comment{Az új osztályok felsorolása. Az régi osztályok közül azoknak a felsorolása, ahol változás volt. Ezek esetén csak a változásokat kell leírni.}

\subsection{Új osztályok}
\subsubsection{Osztály1}
\begin{class-template-responsibility}
    Felelősség leírása
\end{class-template-responsibility}
\begin{class-template-interface}
    Megvalósított interfészek felsorolása
\end{class-template-interface}
\begin{class-template-baseclass}
    Ős-Ősosztály \baseclass Ősosztály... 
\end{class-template-baseclass}
\begin{class-template-attribute}
    \classitem{+A [0..*]}{adattag A}
\end{class-template-attribute}
\comment{Milyen publikus, protected és privát  metódusokkal rendelkezik. Metódusonként precíz leírás, ha szükséges, activity diagram is a metódusban megvalósítandó algoritmusról. Minden olyan metódusnak szerepelnie kell, amelyiket az osztály megvalósít vagy felüldefiniál.}
\begin{class-template-method}
    \classitem{+B(A a) : void}{metódus B}
\end{class-template-method}

\subsection{Megváltozott osztályok}
\subsubsection{Osztály2}
\begin{class-template-responsibility}
    Felelősség leírása
\end{class-template-responsibility}
\begin{class-template-interface}
    Megvalósított interfészek felsorolása
\end{class-template-interface}
\begin{class-template-baseclass}
    Ős-Ősosztály \baseclass Ősosztály... 
\end{class-template-baseclass}
\begin{class-template-attribute}
    \classitem{+A [0..*]}{adattag A}
\end{class-template-attribute}
\comment{Milyen publikus, protected és privát  metódusokkal rendelkezik. Metódusonként precíz leírás, ha szükséges, activity diagram is a metódusban megvalósítandó algoritmusról. Minden olyan metódusnak szerepelnie kell, amelyiket az osztály megvalósít vagy felüldefiniál.}
\begin{class-template-method}
    \classitem{+B(A a) : void}{metódus B}
\end{class-template-method}

\section{Kapcsolat az alkalmazói rendszerrel}
\comment{Szekvencia-diagramokon ábrázolni kell a grafikus rendszer működését. Konzisztens kell legyen az előző alfejezetekkel. Minden metódus, ami ott szerepel, fel kell tűnjön valamelyik szekvenciában. Minden metódusnak, ami szekvenciában szerepel, szereplnie kell a valamelyik osztálydiagramon.}




\clearpage
\section{Napló}

\comment{ 
    A napló tartalmazza az előző beadás óta eltelt időszak történéseit időrendben. A naplóból egyértelműen ki kell derülnie, hogy az egyes anyagrészeket ki és mennyi idő alatt készítette.
    
    A napló bejegyzésekből áll. Minden bejegyzésnek tartalmaznia kell: 
    \begin{itemize}
        \item a történés kezdetének időpontját, nap-óra pontossággal
        \item történés időtartamát, óra felbontással
        \item a szereplő(k) nevét (Kérjük a szereplők VEZETÉKNEVÉT használni)
        \item a tevékenység leírását.
    \end{itemize}
    
    Amennyiben a tevékenységben több szereplő vesz részt, akkor aza tevékenység csak értekezlet lehet, amelynek az eredményei DÖNTÉSEK. A döntéseket precízen meg kell szövegezni (Pl.: Az X objektum Y és Z metódusainak kódját W készíti el Q határidőre).Ha a bejegyzés egyetlen személyhez kötődik, akkor meg kell adni, hogy a tevékenység milyen dologra irányul. A dolog a feladat kapcsán elkészítendő termék, amelynek a (esetleg korábban) beadott anyagban megtalálhatónak kell lenni.A naplóbejegyzés felbontásának egysége szöveges, rajzos anyag esetében az ábra, diagram, vagy kb. fél-egy oldalnyi szöveg. Kódban az egység a metódus. (Pl.: A 3. ábrán látható szekvencia-diagram kidolgozása, vagy az X objektum Y és Z metódusainak kódolása és belövése.
}

\begin{naplo}
    \naplotag{feb. 18. 16h }{ 1 óra }{ Csapat }{ Értekezlet  \newline Döntés: Segédeszközök kiválasztása (git, trello, drive)}
\end{naplo}

\setcounter{chapter}{12}
\chapter{Grafikus felület specifikációja}

\section{Fordítási és futtatási útmutató}
\comment{A feltöltött program fordításával és futtatásával kapcsolatos útmutatás. Ennek tartalmaznia kell leltárszerűen az egyes fájlok pontos nevét, méretét byte-ban, keletkezési idejét, valamint azt, hogy a fájlban mi került megvalósításra.}

\subsection{Fájllista}

\begin{fajllista}
    \fajl
    {Demo.java} % Kezdet
    {353 byte} % Méret
    {2020.03.26~21:05~} % Időpont
    {Demo programosztály} % Leírás
\end{fajllista}

\subsection{Fordítás}
\comment{A fenti listában szereplő forrásfájlokból milyen műveletekkel lehet a bináris, futtatható kódot előállítani. Az előállításhoz csak a 2. Követelmények c. dokumentumban leírt környezetet szabad előírni.}

\begin{verbatim}
javac -d bin *.java
\end{verbatim}

\subsection{Futtatás}
\comment{A futtatható kód elindításával kapcsolatos teendők leírása. Az indításhoz csak a 2. Követelmények c. dokumentumban leírt környezetet szabad előírni.}

\begin{verbatim}
cd bin
java Main.java
\end{verbatim}

\clearpage
\section{Értékelés}
\comment{A projekt kezdete óta az értékelésig eltelt időben tagokra bontva, százalékban.}
\begin{ertekeles}
    \ertekelestag{Teszt János}{ ?????? }{100\%}
\end{ertekeles}


\clearpage
\section{Napló}

\comment{ 
    A napló tartalmazza az előző beadás óta eltelt időszak történéseit időrendben. A naplóból egyértelműen ki kell derülnie, hogy az egyes anyagrészeket ki és mennyi idő alatt készítette.
    
    A napló bejegyzésekből áll. Minden bejegyzésnek tartalmaznia kell: 
    \begin{itemize}
        \item a történés kezdetének időpontját, nap-óra pontossággal
        \item történés időtartamát, óra felbontással
        \item a szereplő(k) nevét (Kérjük a szereplők VEZETÉKNEVÉT használni)
        \item a tevékenység leírását.
    \end{itemize}
    
    Amennyiben a tevékenységben több szereplő vesz részt, akkor aza tevékenység csak értekezlet lehet, amelynek az eredményei DÖNTÉSEK. A döntéseket precízen meg kell szövegezni (Pl.: Az X objektum Y és Z metódusainak kódját W készíti el Q határidőre).Ha a bejegyzés egyetlen személyhez kötődik, akkor meg kell adni, hogy a tevékenység milyen dologra irányul. A dolog a feladat kapcsán elkészítendő termék, amelynek a (esetleg korábban) beadott anyagban megtalálhatónak kell lenni.A naplóbejegyzés felbontásának egysége szöveges, rajzos anyag esetében az ábra, diagram, vagy kb. fél-egy oldalnyi szöveg. Kódban az egység a metódus. (Pl.: A 3. ábrán látható szekvencia-diagram kidolgozása, vagy az X objektum Y és Z metódusainak kódolása és belövése.
}

\begin{naplo}
    \naplotag{feb. 18. 16h }{ 1 óra }{ Csapat }{ Értekezlet  \newline Döntés: Segédeszközök kiválasztása (git, trello, drive)}
\end{naplo}

\chapter{Összefoglalás}
\section{A projektre fordított összes munkaidő}
\comment{Személyenként és összesen}
\begin{ertekelesOra}
    \ertekelestag{TAG1}{NEPTUN1}{ÓRA1}
    \ertekelestag{TAG2}{NEPTUN2}{ÓRA2}
    \ertekelestag{TAG3}{NEPTUN3}{ÓRA3}
    \ertekelestag{TAG4}{NEPTUN4}{ÓRA4}
    \ertekelestag{TAG5}{NEPTUN5}{ÓRA5}
    \ertekelestag{Csapatgyűlés}{ - }{ - }
\end{ertekelesOra}

\section{A feltöltött programok forrássorainak megoszlása}
\begin{ertekelesKod}
    \ertekelestagk{Szkeleton}{??}
    \ertekelestagk{Prototípus}{??}
    \ertekelestagk{Grafikus változat}{??}
    \ertekelestagk{Összesen}{??}
\end{ertekelesKod}

\section{Projekt összegzés}
\comment{A projekt tapasztalatait összegző részben a csapatoknak a projektről kialakult véleményét várjuk. A megválaszolandók köre az alábbi}
\begin{itemize}
\item \textit{Mit tanultak a projektből konkrétan és általában?} \newline
    ...
\item \textit{Mi volt a legnehezebb és a legkönnyebb?} \newline
    ...
\item \textit{Összhangban állt-e az idő és a pontszám az elvégzendő feladatokkal?} \newline
    ...
\item \textit{Ha nem, akkor hol okozott ez nehézséget?} \newline
    ...
\item \textit{Milyen változtatási javaslatuk van?} \newline
    ...
\item \textit{Milyen feladatot ajánlanának a projektre?} \newline
    ...
\end{itemize}




\end{document}